\subsection{Higher dimension fields above the inert quadratic field}
Let $R$ be either $\texttt{R6\_9\_7018}$ or $\texttt{R6\_9\_8785}$.
Set $(c,T)=(195,4405)$ in the first case and $(65,4325)$ in the second.
The full tensors have reduction $\mathbb F_4^3$, and their original
formal codegrees generate $E=\mathbb Q(\sqrt5)$. Their principal ideals
have intersection $c\sqrt5\,\mathcal O_E$. Consequently the signed
sphericalization dimension field $K$ and dimension $\Delta$ satisfy
\[
 E\subseteq K,\qquad \Delta\in c\sqrt5\,\mathcal O_K,\qquad
 0<\operatorname{Tr}_{K/\mathbb Q}(\Delta)\le T.
\]
The codegree ideal assertion is checked directly: a quadratic codegree
with polynomial $X^2-A X+5a^2$ is $a\sqrt5$ times an integral unit of
norm $-1$. The scales $a$ are $65,15,3$ in the first case and $65,5$
in the second; their least common multiples are the stated $c$.

The field $K$ is totally real and abelian by cyclotomicity of the
sphericalization characters. Its degree is $2,4$, or $6$. The finite
free signed order at two has etale special fibre $\mathbb F_4^3$, hence
is the product of three unramified quadratic integer rings. Thus every
completion of $K$ at two has degree two, and its Frobenius at two has
order two. Its restriction to $E$ is nontrivial.

We use the different of the maximal order, without assuming the signed
order is maximal at odd primes. For the integral signed dimensions $d_i$,
character orthogonality gives
\[
 \sum_i d_i\otimes d_i=(\Delta\otimes1)e_K,
\]
where $e_K$ is the separability idempotent in $K\otimes_{\mathbb Q}K$.
If $a\in\mathfrak D_K^{-1}$, contraction of the second tensor factor
by $x\mapsto\operatorname{Tr}_{K/\mathbb Q}(ax)$ gives
\[
 \Delta a=\sum_i d_i\operatorname{Tr}_{K/\mathbb Q}(ad_i)
 \in\mathcal O_K.
\]
Every trace on the right is integral because each $d_i$ is integral.
Thus $\Delta\in\mathfrak D_K$. No inclusion of the maximal order in
the original signed order is used.

\paragraph{Quartic fields.}
The cyclic quartic case is impossible: the unique element of order two
in $C_4$ acts trivially on its quadratic subfield, whereas Frobenius at
two must act nontrivially on $E$.
In the biquadratic case write $K=E F$, with
$F=\mathbb Q(\sqrt q)$, $q>1$ odd squarefree, $5\nmid q$, and
$q\equiv1\pmod4$. Such a choice is always possible by replacing the
second radicand by its product with $5$ if needed. The congruence follows
because $K$ is unramified at two. The discriminants $5,q$ are coprime,
so $\mathcal O_K=\mathcal O_E\otimes\mathcal O_F$ and
$\mathfrak D_K=\mathfrak D_E\mathfrak D_F$.
Write
\[
 \Delta=c\sqrt5\,\frac{a+b\sqrt5}{2},\qquad
 a,b\in\mathcal O_F,\quad a\equiv b\pmod{2\mathcal O_F}.
\]
Total positivity of $\Delta$ implies total positivity of $b$, and
\[
 0<\operatorname{Tr}_F(b)
 =\frac{\operatorname{Tr}_K(\Delta)}{5c}\le B:=\left\lfloor\frac T{5c}\right\rfloor.
\]
Here $B=4$ or $13$. If $p\mid q$ and $p\nmid c$, the different theorem
forces $b$ into the ramified prime $\mathfrak p$ of $F$ above $p$.
The local trace formula
\[
 \operatorname{Tr}(\mathfrak p^k)
 =p^{\lfloor(k+d)/e\rfloor}\mathbb Z_p
\]
with $e=2,d=1,k=1$ gives $p\mid\operatorname{Tr}_F(b)$.
For completeness the trace formula follows from the definition of the
codifferent: $\operatorname{Tr}(\mathfrak p^k)\subseteq p^m\mathbb Z_p$
if and only if $p^{-m}\mathfrak p^k\subseteq\mathfrak p^{-d}$,
equivalently $em\le k+d$. The use of the codifferent is valid because
$\mathfrak p^k$ is an ideal, so it is closed under multiplication by the
local ring of integers.
Therefore
\[
 q/\gcd(q,c)\mid\operatorname{Tr}_F(b),\qquad q/\gcd(q,c)\le B.
\]
For $c=195$, all primes of $q$ are among $3,13$, since every other
possible odd prime exceeds $4$ after excluding $5$.
For $c=65$, its primes are among $3,7,11,13$, with the product of those
other than $13$ at most $13$. Requiring $q\equiv1\pmod4$ leaves exactly
$q=13$ in both cases. This is a proved conductor bound, not a search cutoff.

For the remaining $K=\mathbb Q(\sqrt5,\sqrt{13})$, its maximal order has
product basis $1,(1+\sqrt5)/2,(1+\sqrt{13})/2,(1+\sqrt5)(1+\sqrt{13})/4$.
Put
\[
 \theta=\frac\Delta{c\sqrt5}
 =\frac{x+y\sqrt5+z\sqrt{13}+w\sqrt{65}}4.
\]
Integrality is equivalent to
$x,y,z,w$ having the same parity and $x-y-z+w\equiv0\pmod4$.
Positivity of all conjugates of $\Delta$ gives
\[
 1\le y\le B,\quad x^2<5y^2,\quad13z^2<5y^2,\quad13w^2<y^2,
 \qquad\operatorname{Tr}_K(\Delta)=5cy.
\]
These strict bounds follow by expressing the coefficients of the totally
positive $\Delta$ as signed averages of its conjugates.
The resulting finite enumeration tests positivity and the $d$-number
condition on the degree-four characteristic polynomial
$P_\Delta=X^4+a_1X^3+a_2X^2+a_3X+a_4$:
$a_4^j\mid a_j^4$ for $j=1,2,3,4$. This condition remains necessary when
$\Delta$ itself lies in a proper subfield of $K$.

The two complementary signed codegrees $f,g$ are a Galois-stable pair,
both greater than one. They lie in $K$, hence are either rational integers
or a quadratic pair in $\mathbb Q(\sqrt5)$, $\mathbb Q(\sqrt{13})$ or
$\mathbb Q(\sqrt{65})$. In lowest terms set
\[
 \frac AB=1-\operatorname{Tr}_K(1/\Delta)>0.
\]
Then $f+g=Ak$, $fg=Bk$ for a positive integer $k$, bounded by
$Ak\le T-\operatorname{Tr}_K(\Delta)$. For a nonrational pair require
$Bk\mid(Ak)^2$. In every case check both roots exceed one and compare
the full degree-six Casimir polynomial modulo two to the original tensor.

For $c=195$ there are seven $d$-number candidates and no complementary
pairs. For $c=65$ there are thirty-five candidates; exactly two permit a
pair before parity is imposed. They have
\[
 \Delta=(1625\pm715\sqrt5)/2,\qquad
 P_\Delta=(X^2-1625X+21125)^2,\quad
 (X-f)(X-g)=X^2-143X+169.
\]
The full polynomial reduces to $(X^2+X+1)^3$, whereas the original is
$(X^2+X+1)(X+1)^4$. Thus every quartic case is excluded.

\paragraph{Sextic fields and the complete conductor bound.}
An abelian degree-six field is cyclic, so $K=E F$ for a cyclic cubic field
$F$. Frobenius at two has order two in $K$; its projection to $F$ is
trivial. Hence two splits completely in $F$.
A primitive order-three Dirichlet character has conductor
\[
 f=9^\epsilon\prod_{p\in S}p,\qquad \epsilon\in\{0,1\},\quad
 p\equiv1\pmod3,
\]
with distinct primes in $S$. To see completeness of this shape, at a prime
$p\ne3$ an order-three character is trivial on the pro-$p$ kernel, so its
conductor exponent is one. At three such a character factors through
$(\mathbb Z/9\mathbb Z)^\times$; thus its exponent is two. No prime $5$
occurs, so the discriminants of $E$ and $F$ are coprime.
Again write $\Delta=c\sqrt5(a+b\sqrt5)/2$. Then $b\in\mathcal O_F$ is
totally positive with $\operatorname{Tr}_F(b)\le B$.
For a ramified prime $p$, let $d_p=2$ for tame $p$ and $d_3=4$ when
$9\mid f$. Since its ramification index is three, the different theorem
forces
\[
 b\in\mathfrak p^{k_p},\qquad k_p=\max(0,d_p-3v_p(c)).
\]
This is checked coefficientwise in the product integral basis; at these
odd primes the factor $2$ in $(a+b\sqrt5)/2$ is a unit. The local trace
formula and ideal norm give
\[
 \prod_{p\mid f}p^{\lfloor(k_p+d_p)/3\rfloor}
     \mid\operatorname{Tr}_F(b),\qquad
 \operatorname{Norm}_F(b)\ge\prod_{p\mid f}p^{k_p}.
\]
These divisibilities bound all prime factors and all their simultaneous
occurrences, with no arbitrary conductor cutoff.

For $c=195$, a tame prime other than $13$ contributes itself to the trace
multiple and is at least $7>B=4$. Thus the only possible blocks are $9$
and $13$, yielding conductors $9,13,117$. At conductor $9$ or $13$, two
is inert: modulo nine its image is not $\pm1$, and
$2^4\not\equiv1\pmod{13}$. The only unrejected conductor is $117$.
Here $k_3=1$, so $3\mid\operatorname{Tr}(b)\le4$, giving trace three,
and $\operatorname{Norm}(b)\ge3$. AM--GM gives
$\operatorname{Norm}(b)\le(3/3)^3=1$, a contradiction.

For $c=65$, the only possible tame blocks are $7,13$, and the wild block
is $9$. Blocks $7$ and $9$ together require $63\mid\operatorname{Tr}(b)$,
which is impossible for trace at most $13$. In particular conductors $63$
and $819$ are excluded by the trace bound. All other primes are excluded
because they are at least $19>13$ and would divide that trace.
The complete remaining list is $7,9,13,91,117$. Each single-block field
has two inert: additionally $2^2\not\equiv1\pmod7$.
At conductor $91$, $k_7=2$ and the trace is a positive multiple of seven
at most thirteen, hence exactly seven. Its norm is at least $49$, whereas
AM--GM bounds it by $(7/3)^3<49$. At conductor $117$, $k_3=4$ and the trace
is a positive multiple of nine at most thirteen, hence exactly nine.
Its norm is at least $81$, whereas AM--GM bounds it by $27$.
Every sextic field is therefore excluded.

Consequently both tensors force $K=\mathbb Q(\sqrt5)$ under arbitrary
complex categorification. The exact replay is
\texttt{higher\_fields.py}, and its complete tensor-bound output is
\texttt{certificates.json}. The statement concerns the
signed dimension-character field and does not assume pivotality.

The number-theoretic conductor--discriminant input is Milne,
\emph{Class Field Theory}, Chapter V, Theorem 3.27,
\url{https://www.jmilne.org/math/CourseNotes/CFT.pdf}.
For a cyclic cubic field both nontrivial characters have conductor $f$,
so its discriminant is $f^2$. A ramified prime is totally ramified;
therefore its different exponent is $2$ at a tame conductor prime and
$4$ at the conductor block $9$, as used above.
